\begin{notechapter}
\chapter{Feedforward Neural Networks}\label{ch:feedforward-networks}
\chaptersubtitle{Composing affine maps and nonlinearities, and fitting them by gradient descent}

\begin{chapterintro}
A feedforward neural network is a composition of simple maps: an affine map, a fixed
nonlinearity applied coordinate by coordinate, another affine map, and so on. This chapter
sets up the model and the notation used in the whole part, and the training problem that the
next chapters solve: computing gradients (\cref{ch:backpropagation,ch:computational-graphs}),
following them (\cref{ch:optimizers}), and keeping the process stable
(\cref{ch:activation-functions,ch:normalization,ch:dropout}). You will learn
\begin{itemize}[nosep]
  \item why stacking linear layers gains nothing, and why one nonlinearity changes everything,
        with the example of XOR;
  \item the multilayer perceptron, its parameters, and the column and row conventions;
  \item what the universal approximation theorem says, and what it does not;
  \item the training problem: empirical risk, gradient descent, stochastic and mini-batch
        gradients, epochs.
\end{itemize}
\prerequisites{empirical risk minimization (\cref{ch:erm}); losses as negative
log-likelihoods (\cref{ch:output-distributions}); the perceptron (\cref{ch:perceptron}).}
\end{chapterintro}

\section{From linear models to networks}\label{feedforward-networks:sec:linear}

A linear model predicts with an affine function of the input, $f(x) = w^\top x + b$ for a scalar
output, or $f(x) = Wx + b$ for a vector output. Composing such maps does not enlarge the model:
\[
  W_2(W_1x + b_1) + b_2 = (W_2W_1)x + (W_2b_1 + b_2)
\]
is again affine. A stack of linear layers is a linear model with a factorized matrix (and,
if a middle layer is narrow, a matrix of lower rank: \cref{feedforward-networks:exr:5}).

\paragraph{XOR.} The four points $(0,0)$ and $(1,1)$ with label $0$, $(0,1)$ and $(1,0)$ with
label $1$, cannot be separated by a line. Suppose $w^\top x + b$ were positive on the second pair
and negative on the first. Adding the two inequalities of each pair,
\[
  (w_1 + b) + (w_2 + b) > 0\qquad\text{and}\qquad b + (w_1 + w_2 + b) < 0,
\]
and the two left-hand sides are both equal to $w_1 + w_2 + 2b$: a contradiction. No linear
classifier, and in particular no perceptron (\cref{ch:perceptron}), solves XOR.

\paragraph{One nonlinearity is enough.} Insert a coordinatewise nonlinearity
$\phi$ between two affine maps:
\[
  f(x) = w^\top\phi(Wx + c) + b .
\]
With $\phi = \mathrm{ReLU}$, $\mathrm{ReLU}(t) = \max(0, t)$, the choice (Goodfellow, Bengio and Courville,
\emph{Deep Learning}, MIT Press, 2016, §6.1)
\[
  W = \begin{bmatrix}1&1\\1&1\end{bmatrix},\qquad c = \begin{bmatrix}0\\-1\end{bmatrix},\qquad
  w = \begin{bmatrix}1\\-2\end{bmatrix},\qquad b = 0
\]
computes XOR exactly:

\begin{center}
\begin{tabular}{cccc}
\toprule
$x$ & $Wx + c$ & $h = \mathrm{ReLU}(Wx + c)$ & $f(x) = h_1 - 2h_2$ \\
\midrule
$(0,0)$ & $(0,-1)$ & $(0,0)$ & $0$ \\
$(0,1)$ & $(1,0)$ & $(1,0)$ & $1$ \\
$(1,0)$ & $(1,0)$ & $(1,0)$ & $1$ \\
$(1,1)$ & $(2,1)$ & $(2,1)$ & $0$ \\
\bottomrule
\end{tabular}
\end{center}

The hidden layer has moved the points: $(0,1)$ and $(1,0)$ land on the same point $(1,0)$, and
the three points $(0,0)$, $(1,0)$, $(2,1)$ of the hidden space can be separated by a line
(\cref{feedforward-networks:fig:xor}). The network has \emph{learned a representation} in
which the problem is linear. In this example the representation was designed by hand;
training finds such representations from data.

\begin{figure}[ht]
\centering
\begin{tikzpicture}
\begin{axis}[name=left, width=5.6cm, height=5.6cm, xmin=-0.4, xmax=1.4, ymin=-0.4, ymax=1.4,
  xlabel={$x_1$}, ylabel={$x_2$}, title={input space}, xtick={0,1}, ytick={0,1}, axis on top]
  \addplot[only marks, mark=*, mark size=3pt, bookaccent] coordinates {(0,0) (1,1)};
  \addplot[only marks, mark=square*, mark size=3pt, bookexample] coordinates {(0,1) (1,0)};
\end{axis}
\begin{axis}[at={(left.east)}, anchor=west, xshift=1.6cm, width=5.6cm, height=5.6cm,
  xmin=-0.4, xmax=2.4, ymin=-0.4, ymax=1.4, xlabel={$h_1$}, ylabel={$h_2$},
  title={hidden space}, xtick={0,1,2}, ytick={0,1}, axis on top]
  \addplot[very thick, bookdef, dashed, domain=-0.4:2.4] {(x-0.5)/2};
  \addplot[only marks, mark=*, mark size=3pt, bookaccent] coordinates {(0,0) (2,1)};
  \addplot[only marks, mark=square*, mark size=3pt, bookexample] coordinates {(1,0)};
  \node[font=\scriptsize, anchor=north] at (axis cs:1,-0.05) {$\times2$};
  \node[font=\scriptsize, bookdef, anchor=west] at (axis cs:-0.3,1.2) {line $h_1 - 2h_2 = \frac12$};
\end{axis}
\end{tikzpicture}
\caption{XOR before and after the hidden layer. Circles have label $0$, squares label $1$.
The two squares are sent to the same hidden point $(1,0)$; in the hidden space the
line $h_1 - 2h_2 = \frac12$ separates the classes.}
\label{feedforward-networks:fig:xor}
\end{figure}

\section{The multilayer perceptron}\label{feedforward-networks:sec:mlp}

\begin{definition}[Multilayer perceptron]\label{feedforward-networks:def:mlp}
Let $n_0 = d$ be the input dimension, $n_1,\dots,n_L$ the layer widths, and $\phi$ a function
$\R\to\R$ (the \emph{activation}), applied coordinatewise to vectors. A multilayer perceptron
(MLP) with $L$ layers computes, from $a^{(0)} = x\in\R^d$,
\[
  z^{(l)} = W^{(l)}a^{(l-1)} + b^{(l)},\qquad a^{(l)} = \phi\bigl(z^{(l)}\bigr),\qquad l = 1,\dots,L,
\]
with $W^{(l)}\in\R^{n_l\times n_{l-1}}$ and $b^{(l)}\in\R^{n_l}$. The output is $\hat y = a^{(L)}$, where the last
activation is usually replaced by the identity or by an output unit chosen with the loss
(sigmoid, softmax: \cref{ch:sigmoid-softmax}). $z^{(l)}$ is the \emph{pre-activation} and $a^{(l)}$ the
\emph{activation} of layer $l$; layers $1,\dots,L - 1$ are \emph{hidden}.
\end{definition}

The parameters are $\theta = (W^{(1)}, b^{(1)},\dots,W^{(L)}, b^{(L)})$, and their number is
$\sum_{l=1}^L(n_{l-1} + 1)n_l$. A network $784\to128\to64\to10$ for $28\times28$ images has
$100\,480 + 8\,256 + 650 = 109\,386$ parameters. The names \emph{fully connected network} and,
when $L$ is large, \emph{deep network} are also used. Each coordinate $a^{(l)}_j = \phi(w_j^\top a^{(l-1)} + b_j)$,
with $w_j^\top$ the $j$-th row of $W^{(l)}$, is a \emph{unit} or \emph{neuron}: a linear model followed by
$\phi$. With $\phi = \sign$ it is the perceptron of \cref{ch:perceptron}, which gave the model its
name.

\begin{figure}[ht]
\centering
\begin{tikzpicture}[x=3.6cm, y=1.1cm, >=Stealth,
  unit/.style={circle, draw, thick, minimum size=7mm, inner sep=0pt, fill=white},
  uin/.style={unit, fill=gray!10}, uhid/.style={unit, fill=bookaccent!12}, uout/.style={unit, fill=bookexample!15}]
  \node[uin] (x1) at (0,0.5) {$x_1$};
  \node[uin] (x2) at (0,-0.5) {$x_2$};
  \foreach \i/\y in {1/1,2/0,3/-1} \node[uhid] (h\i) at (1,\y) {$a^{(1)}_{\i}$};
  \node[uout] (o) at (2,0) {$\hat y$};
  \foreach \i in {1,2} \foreach \j in {1,2,3} \draw[->, gray] (x\i) -- (h\j);
  \foreach \j in {1,2,3} \draw[->, gray] (h\j) -- (o);
  \draw[->, thick, bookaccent] (x2) -- (h1) node[pos=0.3, above left, font=\scriptsize, bookaccent] {$W^{(1)}_{12}$};
  \node[font=\small] at (0,-1.8) {input $x$};
  \node[font=\small] at (1,-1.8) {hidden layer $a^{(1)}$};
  \node[font=\small] at (2,-1.8) {output $\hat y$};
\end{tikzpicture}
\caption{An MLP $2\to3\to1$, computing $a^{(1)} = \phi(W^{(1)}x + b^{(1)})$ and $\hat y = \sigma(W^{(2)}a^{(1)} + b^{(2)})$: the network of
the worked example of \cref{ch:backpropagation}.
The weight $W^{(1)}_{jk}$ connects input $k$ to hidden unit $j$: row $j$ of $W^{(1)}$ holds the
incoming weights of unit $j$.}
\label{feedforward-networks:fig:mlp}
\end{figure}

\begin{remark}[Column and row conventions]\label{feedforward-networks:rem:rows}
This part writes vectors as columns and a layer as $z = Wa + b$, so that $W_{jk}$ is the weight
from unit $k$ to unit $j$. Software processes a mini-batch of $B$ examples stored as the rows of
a matrix $A\in\R^{B\times n_{l-1}}$ and computes $Z = AW^\top + \mathbf 1b^\top$, whose row $i$ is $(Wa_i + b)^\top$.
PyTorch's \texttt{nn.Linear(n\_in, n\_out)} stores \texttt{weight} with shape
$(n_{\mathrm{out}}, n_{\mathrm{in}})$, the same $W$ as here, and computes
\texttt{x @ weight.T + bias}. The two conventions describe the same map;
\cref{backpropagation:sec:batch} gives the gradients in the row form.
\end{remark}

\section{What networks can represent}\label{feedforward-networks:sec:uat}

\begin{theorem}[Universal approximation]\label{feedforward-networks:thm:uat}
Let $\phi:\R\to\R$ be continuous. The functions $x\mapsto\sum_{j=1}^Nv_j\phi(w_j^\top x + c_j)$, for all $N$ and
all $v_j, c_j\in\R$, $w_j\in\R^d$, are dense in $C(K)$ for the uniform norm, for every compact $K\subset\R^d$, if and
only if $\phi$ is not a polynomial.
\end{theorem}

This form is due to Leshno, Lin, Pinkus and Schocken (\emph{Multilayer Feedforward Networks
with a Nonpolynomial Activation Function Can Approximate Any Function}, Neural Networks,
1993); earlier versions are Cybenko (\emph{Approximation by Superpositions of a Sigmoidal
Function}, Mathematics of Control, Signals and Systems, 1989) for sigmoidal $\phi$ and Hornik
(\emph{Approximation Capabilities of Multilayer Feedforward Networks}, Neural Networks,
1991). The ``only if'' is easy: if $\phi$ is a polynomial of degree $k$, every such sum is a
polynomial of degree at most $k$, a closed proper subspace of $C(K)$.

\begin{pitfall}[What the theorem does not say]
It is an existence result for one hidden layer. It gives no bound on the width $N$ needed
(which can grow exponentially with $d$ for some functions), no algorithm to find the
weights, and no statement about generalization from finitely many examples. Its value is
negative: the choice of a (non-polynomial) activation does not limit what can be
represented, so the questions worth asking are about efficiency, optimization and
generalization.
\end{pitfall}

\paragraph{Depth.} Depth can be exponentially more efficient than width. The tent map
\[
  t(x) = 2\,\mathrm{ReLU}(x) - 4\,\mathrm{ReLU}\bigl(x - \tfrac12\bigr)
\]
maps $[0,1]$ onto itself, rising then falling. Composing
it $k$ times gives a sawtooth with $2^{k-1}$ teeth, that is $2^k$ linear pieces
(\cref{feedforward-networks:exr:3}), with a network of depth $k$ and width $2$. A network
with one hidden layer of $N$ ReLU units is piecewise linear in one variable with at most
$N + 1$ pieces, so it needs $N\geq2^k - 1$ units for the same function. Telgarsky
(\emph{Benefits of Depth in Neural Networks}, COLT, 2016) turned this into a
separation theorem.

\begin{remark}
Kolmogorov's superposition theorem (1957) states that every continuous function on $[0,1]^d$ can
be written exactly with sums and continuous functions of one variable, with a fixed number
$2d + 1$ of outer terms. It is sometimes quoted as an early ``representation theorem'' for
networks, but its inner functions are highly irregular and are not of the form
$\phi(w^\top x + c)$ with a fixed $\phi$.
\end{remark}

\section{Training}\label{feedforward-networks:sec:training}

\paragraph{The objective.} With a loss $\ell(\hat y, y)$ and training data $(x_1,y_1),\dots,(x_n,y_n)$, training
minimizes the empirical risk (\cref{ch:erm})
\[
  \ERisk(\theta) = \frac1n\sum_{i=1}^n\ell\bigl(f_\theta(x_i), y_i\bigr),
\]
a surrogate for the population risk $\Risk(\theta) = \E[\ell(f_\theta(X), Y)]$. The loss is chosen as the
negative log-likelihood of an output distribution (\cref{ch:output-distributions}): the
squared error $\frac12\norm{\hat y - y}^2$ for a Gaussian output, the cross-entropy $-\sum_ky_k\log\hat y_k$ with a
softmax output for classification (\cref{ch:sigmoid-softmax}). Because of the
nonlinearities, $\ERisk$ is not convex in $\theta$; it is minimized by gradient methods, which only
need $\nabla\ERisk(\theta)$, computed by backpropagation (\cref{ch:backpropagation}).

\paragraph{Gradient descent.} Full-batch gradient descent iterates
\[
  \theta_{t+1} = \theta_t - \eta\,\nabla\ERisk(\theta_t),
\]
with a \emph{learning rate} $\eta>0$. Each step costs a pass over the $n$ examples.

\paragraph{Stochastic and mini-batch gradients.} Since $\nabla\ERisk = \frac1n\sum_i\nabla\ell_i$ is an average, it can be
estimated from a sample. Draw a \emph{mini-batch} $\mathcal B$ of $B$ indices uniformly and use
\[
  g_{\mathcal B}(\theta) = \frac1B\sum_{i\in\mathcal B}\nabla_\theta\ell\bigl(f_\theta(x_i), y_i\bigr),\qquad
  \theta_{t+1} = \theta_t - \eta\,g_{\mathcal B_t}(\theta_t) .
\]
$B = 1$ is \emph{stochastic gradient descent} (SGD) in the strict sense, $B = n$ is gradient
descent, and the intermediate case, mini-batch SGD, is what is used in practice and is
also simply called SGD.

\begin{proposition}[Mini-batch gradients are unbiased]\label{feedforward-networks:prop:minibatch}
If the $B$ indices are drawn independently and uniformly from $\{1,\dots,n\}$, then
$\E[g_{\mathcal B}] = \nabla\ERisk$ and $\operatorname{Cov}(g_{\mathcal B}) = \frac1B\Sigma$, where
$\Sigma = \frac1n\sum_i(\nabla\ell_i - \nabla\ERisk)(\nabla\ell_i - \nabla\ERisk)^\top$ is the covariance of the per-example gradients.
Without replacement, the covariance is multiplied by $\frac{n - B}{n - 1}$.
\end{proposition}

\begin{proof}
Each term $\nabla\ell_{I}$ with $I$ uniform has mean $\frac1n\sum_i\nabla\ell_i = \nabla\ERisk$ and covariance $\Sigma$; the average of $B$
independent copies has the same mean and covariance $\Sigma/B$. The finite-population factor for
sampling without replacement is the classical one of survey sampling
(\cref{feedforward-networks:exr:4}).
\end{proof}

The noise of the estimate decreases like $1/\sqrt B$ while the cost of a step grows like $B$:
small batches give many cheap, noisy steps. The noise is not only a nuisance; it is thought
to help escape poor regions of the loss surface, and it interacts with the learning rate
(the ratio $\eta/B$ matters in many analyses).

\paragraph{Epochs, validation, schedules.} One \emph{epoch} is a pass over the training set,
usually by shuffling the indices once and cutting them into consecutive batches (sampling
without replacement). The learning rate, the batch size and the number of epochs are
hyperparameters chosen on a validation set (\cref{ch:measuring-performance}); the learning
rate is often varied during training by a \emph{schedule} (\cref{optimizers:sec:schedules}).

\paragraph{In code.} In PyTorch the gradient is computed by automatic differentiation
(\cref{ch:computational-graphs}) and the update by an optimizer object:
\begin{lstlisting}[language=Python]
import torch
from torch import nn

model = nn.Sequential(nn.Linear(784, 128), nn.ReLU(),
                      nn.Linear(128, 64), nn.ReLU(),
                      nn.Linear(64, 10))          # logits; softmax is in the loss
loss_fn = nn.CrossEntropyLoss()
opt = torch.optim.SGD(model.parameters(), lr=0.1)

for epoch in range(10):
    model.train()
    for xb, yb in train_loader:                   # shuffled mini-batches
        opt.zero_grad()                           # gradients accumulate otherwise
        loss = loss_fn(model(xb.flatten(1)), yb)  # forward pass
        loss.backward()                           # backpropagation
        opt.step()                                # theta <- theta - lr * g
\end{lstlisting}
\texttt{train\_loader} is a \texttt{DataLoader} over the training set with
\texttt{shuffle=True}. \texttt{model.train()} matters for layers that behave differently at
training and test time, dropout and batch normalization (\cref{ch:dropout,ch:normalization});
\texttt{model.eval()} switches them to inference mode.

\begin{chaptersummary}
\begin{itemize}[nosep]
  \item Compositions of affine maps are affine; XOR is not linearly separable. One hidden layer
        with a nonlinearity solves it, by mapping the inputs to a representation where the
        classes are separable.
  \item An MLP alternates $z^{(l)} = W^{(l)}a^{(l-1)} + b^{(l)}$ and $a^{(l)} = \phi(z^{(l)})$; it has $\sum_l(n_{l-1} + 1)n_l$ parameters.
        Software uses rows: $Z = AW^\top + \mathbf 1b^\top$.
  \item One hidden layer with a continuous non-polynomial activation can approximate any
        continuous function on a compact set, with no bound on the width; depth can be
        exponentially more economical.
  \item Training minimizes the empirical risk by (mini-batch) stochastic gradient descent; mini-batch
        gradients are unbiased with covariance $\Sigma/B$.
\end{itemize}
\end{chaptersummary}

\section*{Exercises}
\addcontentsline{toc}{section}{Exercises}

\begin{exercise}\label{feedforward-networks:exr:1}
Count the parameters of an MLP $3\to5\to5\to2$. How many does a linear model $3\to2$ have,
and which linear models can the MLP represent if $\phi$ is the identity?
\solutionlink{feedforward-networks:exr:1}
\end{exercise}

\begin{exercise}\label{feedforward-networks:exr:2}
Find a network $w^\top\mathrm{ReLU}(Wx + c) + b$ with two hidden units that computes the ``equivalence''
function, $1$ on $(0,0)$ and $(1,1)$ and $0$ on $(0,1)$ and $(1,0)$.
\solutionlink{feedforward-networks:exr:2}
\end{exercise}

\begin{exercise}\label{feedforward-networks:exr:3}
Let $t(x) = 2\,\mathrm{ReLU}(x) - 4\,\mathrm{ReLU}(x - \frac12)$. Show that $t(x) = 2x$ on $[0,\frac12]$ and $t(x) = 2 - 2x$ on $[\frac12,1]$, and
that $t^{\circ k} = t\circ\dots\circ t$ ($k$ times) is piecewise linear on $[0,1]$ with $2^k$ pieces.
\solutionlink{feedforward-networks:exr:3}
\end{exercise}

\begin{exercise}\label{feedforward-networks:exr:4}
Let $u_1,\dots,u_n\in\R$ with mean $\bar u$ and $s^2 = \frac1n\sum_i(u_i - \bar u)^2$. Show that the mean of a sample of $B$
values drawn without replacement has variance $\frac{s^2}B\frac{n - B}{n - 1}$. What happens for $B = n$?
\solutionlink{feedforward-networks:exr:4}
\end{exercise}

\begin{exercise}\label{feedforward-networks:exr:5}
An MLP $d\to k\to m$ with identity activation computes $x\mapsto W_2W_1x + c$. Show that the set of
matrices $W_2W_1$ is exactly the set of $m\times d$ matrices of rank at most $\min(k, d, m)$. Which known
method does least-squares training of this network with $k<\min(d,m)$ and $m = d$, $y = x$
correspond to?
\solutionlink{feedforward-networks:exr:5}
\end{exercise}
\end{notechapter}
